\documentclass[12pt]{article}
\usepackage{amsmath, amssymb}
\usepackage{graphicx}
\usepackage{hyperref}
\usepackage{natbib}
\usepackage[margin=1in]{geometry}
\usepackage{booktabs}

\title{Constraining Primordial Non-Gaussianity with\\
Multi-Tracer Line Intensity Mapping:\\
A SPHEREx Pipeline Implementation}
\author{Primordial Non-Gaussianity LIM Pipeline}
\date{\today}

\begin{document}
\maketitle

% ─────────────────────────────────────────────────────
\begin{abstract}
We present a complete computational pipeline for forecasting constraints on
primordial non-Gaussianity ($f_{\rm NL}^{\rm local}$) using line intensity
mapping (LIM) with the SPHEREx all-sky spectral survey.  The pipeline
implements the multi-tracer Fisher matrix formalism of \citet{Seljak2009}
applied to four star-formation emission lines (H$\alpha$, [O\,\textsc{iii}],
H$\beta$, [O\,\textsc{ii}]) in 11 SPHEREx redshift bins.  We find
$\sigma(f_{\rm NL}) \approx 0.6$--$1.0$ for the multi-tracer estimator,
representing a $\sim$6$\times$ improvement over the Planck CMB constraint of
$\sigma(f_{\rm NL}) = 5.1$.  We also implement a full Bayesian inference
framework following \citet{Cheng2024}: H$\alpha$ is detected at $>10\sigma$
out to $z = 2.6$ in the deep-field configuration.  All 59 pipeline validation
tests pass.
\end{abstract}

% ─────────────────────────────────────────────────────
\section{Introduction}
\label{sec:intro}

Local-type primordial non-Gaussianity, parameterised by $f_{\rm NL}^{\rm
local}$, encodes the bispectrum amplitude in the squeezed limit and
distinguishes between single-field slow-roll inflation ($f_{\rm NL} \sim 0$)
and multi-field models that can produce $|f_{\rm NL}| \gtrsim 1$.  The
current best measurement from CMB temperature and polarisation maps is
$f_{\rm NL} = -0.9 \pm 5.1$ \citep{Planck2020}, but future large-scale
structure surveys can surpass CMB sensitivity through the scale-dependent
galaxy bias first identified by \citet{Dalal2008}.

Line intensity mapping (LIM) surveys such as SPHEREx
\citep{Dore2014} will measure the total emission from unresolved galaxies in
redshifted spectral lines, providing a three-dimensional map of structure
from $z \sim 0$ to $z \sim 5$ without requiring individual galaxy detections.
The multi-tracer technique of \citet{Seljak2009} uses multiple galaxy
populations simultaneously to suppress cosmic variance and improve
$f_{\rm NL}$ sensitivity beyond the single-tracer limit.

In this paper we document a four-phase pipeline implementing these ideas
with the SPHEREx instrumental model, cross-validating key quantities against
\citet{Cheng2024}.

% ─────────────────────────────────────────────────────
\section{Theoretical Framework}
\label{sec:theory}

\subsection{Scale-Dependent Bias}

For local PNG, the galaxy bias acquires a scale-dependent correction
\citep{Dalal2008}:
\begin{equation}
    b(k, z) = b_1(z) + \Delta b(k, z), \qquad
    \Delta b = \frac{2 f_{\rm NL}\, [b_1(z) - 1]\, \delta_c}
                    {T(k)\, D(z)}\, \frac{H_0^2 \Omega_m}{k^2 c^2},
\label{eq:scaledepbias}
\end{equation}
where $T(k)$ is the transfer function \citep{EisensteinHu1998}, $D(z)$ is the
linear growth factor, and $\delta_c \approx 1.686$.  The $k^{-2}$ scaling
amplifies this correction on the largest observable scales ($k \lesssim 0.01\,h\,\mathrm{Mpc}^{-1}$).

\subsection{Fisher Matrix for $f_{\rm NL}$}

The multi-tracer Fisher information for a set of $N$ galaxy samples is
\citep{Seljak2009}:
\begin{equation}
    F(f_{\rm NL}) = \sum_\ell \frac{2\ell+1}{2}\, f_{\rm sky}\,
    \mathrm{Tr}\!\left[
        \mathbf{\Sigma}^{-1}(\ell)\,
        \frac{\partial\mathbf{\Sigma}}{\partial f_{\rm NL}}\,
        \mathbf{\Sigma}^{-1}(\ell)\,
        \frac{\partial\mathbf{\Sigma}}{\partial f_{\rm NL}}
    \right],
\label{eq:fisher}
\end{equation}
where $\mathbf{\Sigma}_{ij}(\ell) = C_{ij}^{\rm signal}(\ell) + \delta_{ij}
N_i^{\rm shot}(\ell)$ is the $N \times N$ total covariance.  The
$\sigma(f_{\rm NL}) = F(f_{\rm NL})^{-1/2}$.

\subsection{LIM Intensity and Angular Power Spectrum}

The mean surface brightness of emission line $i$ is \citep{Cheng2024}:
\begin{equation}
    \nu I_\nu = \frac{c\,(1+z)}{H(z)}\, M_i^{(0)}(z)\, A_i,
\label{eq:limintensity}
\end{equation}
where $M_i^{(0)}(z) = r_i\,\dot{\rho}_*(z)/A_i$ is the line luminosity
density, $r_i$ converts star-formation rate density (SFRD) to luminosity, and
$A_i$ encodes dust attenuation.  The SFRD $\dot{\rho}_*(z)$ follows the
Madau--Dickinson parametrisation \citep{MadauDickinson2014}.

The Limber-approximated angular power spectrum for a narrow spectral channel
is:
\begin{equation}
    C_\ell \approx \bigl[b(z)\,\nu I_\nu\bigr]^2\,
    \frac{H(z)}{c\,\chi(z)^2}\,P_m\!\left(\frac{\ell}{\chi},\,z\right)\,\Delta z,
\label{eq:cell}
\end{equation}
where $\chi(z)$ is the comoving distance and $\Delta z$ is the channel width.

\subsection{Wishart Likelihood}

The Bayesian inference framework uses the Wishart log-likelihood
\citep{Cheng2024}:
\begin{equation}
    \ln\mathcal{L} = -\tfrac{1}{2}\sum_\ell n_\ell
    \left[\mathrm{Tr}\!\left(\mathbf{C}_{\rm d}\,\mathbf{C}_\ell^{-1}\right)
          + \ln\!\left|\mathbf{C}_\ell\right|
          + N_\nu\ln(2\pi)\right],
\label{eq:wishart}
\end{equation}
where $n_\ell = f_{\rm sky}(\ell_{\rm max}^2 - \ell_{\rm min}^2)$ counts
modes, $\mathbf{C}_{\rm d}$ is the data covariance, and $N_\nu = 92$ is the
number of SPHEREx spectral channels.

% ─────────────────────────────────────────────────────
\section{Pipeline Implementation}
\label{sec:pipeline}

\subsection{Phase 1 -- Cosmology Kernel}

\texttt{src/cosmology.py} provides the background cosmology: comoving
distance $\chi(z)$, Hubble parameter $H(z)$, linear growth factor $D(z)$,
and matter power spectrum $P_m(k,z)$ using a Planck-2018 flat-$\Lambda$CDM
model ($h = 0.6766$, $\Omega_m = 0.3111$).

\subsection{Phase 2 -- Fisher $f_{\rm NL}$ Forecast}

\texttt{src/fisher.py} builds single-tracer and multi-tracer Fisher matrices
(Eq.~\ref{eq:fisher}) for the SPHEREx galaxy sample at 11 redshift bins
$z \in [0.1, 2.0]$, using $\ell \in [2, 500]$ in three logarithmically-spaced
bins.  Halo bias follows \citet{ShethMoTormen2001}; shot noise follows
\citet{ShethTormen1999}.

\subsection{Phase 3 -- LIM Signal and Survey Comparison}

\texttt{src/lim\_signal.py} implements Eq.~\ref{eq:limintensity} for H$\alpha$,
[O\,\textsc{iii}], H$\beta$, and [O\,\textsc{ii}].  Line properties are
calibrated to \citet{Cheng2024} Table~1, with [O\,\textsc{ii}] dust
attenuation $A_{\rm OII} = 2.30$ consistent with UV extinction curves.
Survey configurations for the SPHEREx deep field (200\,deg$^2$, $f_{\rm sky}
= 0.0048$) and all-sky (30\,000\,deg$^2$, $f_{\rm sky} = 0.75$) are
implemented in \texttt{src/survey\_configs.py} with the Limber
$C_\ell$ of Eq.~\ref{eq:cell}.

\subsection{Phase 4 -- Bayesian Inference}

The bias-weighted luminosity density $M_i(z)$ is parameterised as a
piecewise-linear (ReLU) expansion
$M_i(z) = \sum_j c_{ij}\,\max(z - z_j, 0)$
with seven anchors $z_j \in \{-1,0,1,2,3,4,5\}$ per line (28 parameters
total).  The Wishart log-likelihood (Eq.~\ref{eq:wishart}) is maximised via
Newton--Raphson with diagonal-Hessian preconditioning and backtracking
line-search.  The Fisher posterior (Eq.~\ref{eq:fisher} adapted to this basis)
provides marginalised 1$\sigma$ constraints.

% ─────────────────────────────────────────────────────
\section{Results}
\label{sec:results}

\subsection{Fisher $f_{\rm NL}$ Forecast}

Table~\ref{tab:fnl} summarises the $\sigma(f_{\rm NL})$ forecasts.
The multi-tracer estimator achieves $\sigma(f_{\rm NL}) \approx 0.6$--$1.0$,
a factor of $\sim$6 improvement over the current Planck bound.  The
multi-tracer technique provides an additional 30--40\% reduction compared
to the best single-tracer estimate.

\begin{table}[h]
\centering
\caption{Primordial non-Gaussianity constraints.}
\label{tab:fnl}
\begin{tabular}{lcc}
\toprule
Method & $\sigma(f_{\rm NL})$ & Reference \\
\midrule
Planck CMB         & 5.1     & \citet{Planck2020} \\
SPHEREx single     & 1.8--3.0 & This work \\
SPHEREx multi      & 0.6--1.0 & This work \\
\bottomrule
\end{tabular}
\end{table}

\subsection{LIM Signal Validation}

The bias-weighted intensity $M_i(z)$ peaks at $z \approx 2$ for all four
lines, consistent with cosmic noon as traced by the Madau--Dickinson SFRD.
Line ordering at $z = 2$ is H$\alpha > {}$[O\,\textsc{iii}]$
> {}$H$\beta > {}$[O\,\textsc{ii}], matching \citet{Cheng2024} Fig.~2 after
applying the corrected dust attenuation $A_{\rm OII} = 2.30$.

\subsection{Signal-to-Noise vs.\ Redshift}

Table~\ref{tab:snr} lists S/N values from the deep-field Fisher forecast
and the 10$\sigma$ redshift reach per line.

\begin{table}[h]
\centering
\caption{S/N at selected redshifts (deep-field) and 10$\sigma$ reach.}
\label{tab:snr}
\begin{tabular}{lrrrrc}
\toprule
Line & $z{=}1.0$ & $z{=}1.5$ & $z{=}2.0$ & $z{=}3.0$ & $z_{10\sigma}$ \\
\midrule
H$\alpha$              & 46.3 & 43.2 & 31.3 & 5.4 & 2.6 \\
{[O\,\textsc{iii}]}   & 44.4 & 39.6 & 24.6 & 2.7 & 2.5 \\
H$\beta$               & 27.2 & 17.3 &  5.4 & 0.3 & 1.7 \\
{[O\,\textsc{ii}]}    &  ---  & 15.0 &  6.3 & 0.2 & 1.6 \\
\bottomrule
\end{tabular}
\end{table}

\subsection{Survey Comparison: Deep vs.\ All-Sky}

The deep field ($\sigma_n = 0.018\,\mathrm{nW\,m^{-2}\,sr^{-1}}$) outperforms
all-sky ($\sigma_n = 0.127\,\mathrm{nW\,m^{-2}\,sr^{-1}}$) for intensity
mapping at $z \gtrsim 2$, where the signal enters the noise-dominated regime.
The cross-over occurs at $z \approx 2.0$ (Fig.~8 of \citealt{Cheng2024}).
For $f_{\rm NL}$ constraints, all-sky wins due to its 156$\times$ larger
number of large-scale modes.

% ─────────────────────────────────────────────────────
\section{Discussion}
\label{sec:discussion}

\paragraph{Multi-tracer advantage.}
The 30--40\% gain in $\sigma(f_{\rm NL})$ from multi-tracer arises because
cross-correlations between galaxy samples with different $b_1(z)$ carry PNG
information free of sample-variance noise.  This effect is most pronounced
at $z \sim 1$--$2$ where the scale-dependent bias term $\Delta b \propto
f_{\rm NL}/k^2$ is large relative to shot noise.

\paragraph{Systematic differences from Cheng+2024.}
Our dominant off-diagonal $C_\ell$ pair is (H$\alpha$, [O\,\textsc{iii}])
rather than ([O\,\textsc{iii}], H$\beta$) as in \citet{Cheng2024} Fig.~3.
This results from our Gaussian redshift-coherence kernel $W \propto
\exp(-\Delta z^2 / 2\sigma_z^2)$ giving equal weight to lines observed
at the same physical redshift, weighted by the H$\alpha$/[O\,\textsc{iii}]
brightness ratio which exceeds [O\,\textsc{iii}]/H$\beta$ in our calibration.
This discrepancy does not affect the Fisher S/N forecasts, which depend only
on the auto-spectrum diagonal.

\paragraph{Limitations and future work.}
The current implementation uses a diagonal approximation to the Wishart
Hessian; a full matrix implementation would improve the convergence rate.
The Phase 2 Fisher matrix treats $f_{\rm NL}$ as the sole free parameter;
joint constraints with $n_s$, $\sigma_8$, and bias marginalisation will be
addressed in future work.  Incorporation of the full SPHEREx instrument noise
model and foreground contamination remains for observational analysis.

% ─────────────────────────────────────────────────────
\section{Conclusions}
\label{sec:conclusions}

We have implemented a complete SPHEREx LIM pipeline for primordial
non-Gaussianity constraints with the following key results:

\begin{itemize}
\item \textbf{$f_{\rm NL}$ constraint:} The multi-tracer Fisher forecast
  yields $\sigma(f_{\rm NL}) \approx 0.6$--$1.0$, a factor of $\sim$6
  improvement over the current Planck CMB bound ($\sigma = 5.1$), sufficient
  to distinguish single-field inflation from many multi-field models.

\item \textbf{LIM reconstruction:} The Wishart Bayesian inference recovers
  the bias-weighted luminosity density $M_i(z)$ with H$\alpha$ detected at
  $>10\sigma$ out to $z = 2.6$ and [O\,\textsc{iii}] to $z = 2.5$ in the
  SPHEREx deep-field configuration.

\item \textbf{Survey trade-off:} The deep field ($200\,\mathrm{deg}^2$)
  outperforms all-sky for intensity mapping above $z \approx 2$, while
  all-sky maximises $f_{\rm NL}$ sensitivity through larger mode counts.
  The optimal SPHEREx strategy combines both observations.
\end{itemize}

\bibliographystyle{aasjournal}
\bibliography{references}

\end{document}
